\section{A Three-Boundary Verification Model}
\label{sec:model}

We model a turn as
\begin{equation}
R=(\tau, a, o, f, S, E, Y, z),
\end{equation}
where $\tau$ is a trace identifier, $a$ an authenticated actor, $o$ an
ephemeral lease-owner token, $f$ a monotonic fencing token, $S$ a
request-scoped clinical state, $E$ the validated evidence records, $Y$ the
private attempt outputs, and $z$ the durable terminal state.  This notation is
a reference contract, not a formal safety theorem.  Table~\ref{tab:status}
separates the target obligation from the present implementation.

\subsection{Claim admissibility}

Let $s$ be a detected medical text segment and
$\operatorname{marker}(s,E)$ indicate that the segment contains an in-range
current-turn evidence marker.  The implemented audit is
\begin{equation}
\operatorname{ClaimAudit}(s,E)=
\begin{cases}
\textsc{record-bound}, & \operatorname{marker}(s,E),\\
\textsc{record-unbound}, & \neg\operatorname{marker}(s,E).
\end{cases}
\end{equation}
For a narrow regex-defined class $D$ of deterministic diagnosis assertions,
the output sanitizer additionally applies
\begin{equation}
\claimcheck(s,E)=
\begin{cases}
\textsc{retain-text}, & s\in D\wedge\operatorname{marker}(s,E),\\
\textsc{rewrite-uncertain}, & s\in D\wedge\neg\operatorname{marker}(s,E),\\
\textsc{audit-only}, & s\notin D.
\end{cases}
\end{equation}
Thus a bound marker means structurally traceable, not semantically supported.
General medication, risk, and clinical statements are recorded but do not all
fail closed.  Before this path, deterministic red-flag patterns can produce
urgent-action guidance without a model call.

\subsection{Run finality}

Public finality is allowed only while both coordination layers agree:
\begin{equation}
\runcommit(R) =
\operatorname{ownsRedis}(o)
\wedge \operatorname{currentPG}(a,\tau,f)
\wedge \operatorname{atomic}(m,q,z),
\end{equation}
where $m$ is the assistant message and $q$ the audit event.  A worker may
stream provisional content, but it cannot publish \texttt{done} until
$(m,q,z)$ commits.  An older worker with a stale $f$ cannot write either a
successful or failed terminal state after takeover.  Attempt events remain
private; only the committed state becomes public.

\subsection{Update admissibility}

For candidate $u$ and baseline $b$, the target update contract is
\begin{align}
\updategate(u,b)={}&
\operatorname{frozen}(u)
\wedge \operatorname{sameRunner}(u,b)
\wedge \operatorname{noPublicRegression}(u,b)\nonumber\\
&\wedge \operatorname{sealedAttestation}(u)
\wedge \operatorname{humanSignature}(u)
\wedge \operatorname{atomicPromotion}(u).
\end{align}
The current concrete runner covers only \texttt{routing.strategy} with four
deterministic cases, one per required slice.  The repository defines signed
envelopes for sealed-gate facts and exact-artifact human approval, plus atomic
Git reference promotion.  It does not implement the service that runs hidden
cases or derives those facts, nor deploy the signer as a separately
authenticated authority.  Those are requirements of \updategate, not achieved
properties of the present prototype.

\begin{table}[t]
\caption{Reference obligations and implementation status.}
\label{tab:status}
\centering
\scriptsize
\begin{tabularx}{\linewidth}{@{}p{0.11\linewidth}X X X@{}}
\toprule
Object & Reference obligation & Present implementation & Not established \\
\midrule
Claim & Admit only claims satisfying an evidence policy & In-range marker
normalization/audit; diagnosis-pattern rewrite; emergency short-circuit &
Entailment and fail-closed handling of all clinical statements \\
Run & Only the current owner atomically exposes one terminal state & Redis
lease; monotonic PostgreSQL fencing; row lock; shared terminal transaction &
Tolerance of datastore or host compromise; exactly-once downstream delivery \\
Update & Same-runner regression gate plus independent evaluation and approval
& Candidate freeze/sandbox; four-case routing runner; signed-envelope and
atomic-promotion libraries & Hidden-case execution and separately deployed
evaluator/approver authorities \\
\bottomrule
\end{tabularx}
\end{table}

\subsection{Threat model and trusted computing base}

Model outputs, retrieved pages, uploaded documents, tools, and candidate
bundles are untrusted.  A runtime worker may crash, retry, or become stale; a
candidate may attempt to inspect evaluator inputs or write outside its bundle.
We trust PostgreSQL transaction and locking semantics, Redis script atomicity
for lease coordination, the controller and Docker host, protected Git
references, and uncompromised signing keys.  The run contract does not survive
database or host compromise.  Update signatures do not provide independence
if a controller can access evaluator or approval keys.  Network partitions
may reduce availability; a stale worker remains rejectable only if the
authoritative database and fencing check are reachable.  These assumptions
define the trusted computing base rather than proving safety under a fully
compromised platform.
